%% ---- wide page ------------------------------------------------------------
%% It widens the by gettint the maximum of the estimation of wideness provided by longtable (2pass)
%% or the provided parameter. If it is no table it uses the provided number directly.

\newlength{\wpcharwd}
\newlength{\wptblwd}
\newlength{\wpmargin}
\newlength{\wpminmargin}    % margin needed to keep the watermark clear of text
\newlength{\wpcontentwd}    % widest table measured so far in the current block
\newlength{\wpscratch}
\newlength{\wpsaved}
\newlength{\wpsavedtw}
\newlength{\wpsavedosm}
\newlength{\wpsavedesm}
\newcounter{wpblock}
\newif\ifwpwidened
\newif\ifwpinblock

\ifdefined\pdfpagewidth
  \let\wppagewidth\pdfpagewidth      % pdfTeX, XeTeX
\else
  \let\wppagewidth\pagewidth         % LuaTeX
\fi

\newcommand{\wptabfont}{\ttfamily\fontsize{7}{8}\selectfont}

\makeatletter

%% Width of a block's widest table, as recorded by the previous run.
\gdef\wpsetwidth#1#2{\global\@namedef{wp@wd@#1}{#2}}

%% Sum of the column widths longtable just settled on == the table's real width.
\def\wp@addentry#1#2{\advance\wpscratch#2\relax}
\def\wp@measure{%
  \ifwpinblock
    \begingroup
      \wpscratch\z@
      \let\LT@entry\wp@addentry
      \LT@save@row
      \ifdim\wpscratch>\wpcontentwd
        \global\wpcontentwd\wpscratch
      \fi
    \endgroup
  \fi}

%% longtable is loaded after this file, so hook it once the preamble is done.
\AtBeginDocument{%
  \@ifundefined{wpstamp}{}{%
    \settowidth{\wpscratch}{\wpstamp}%
    \global\wpminmargin\dimexpr\wpstampinset+\wpscratch+\wpstampgap\relax}%
  \@ifundefined{endlongtable}{}{%
    \let\wp@endlongtable\endlongtable
    \def\endlongtable{\wp@endlongtable\wp@measure}}%
}

\def\wp@rerun{%
  \PackageWarning{widepage}{Wide page \thewpblock\space width has changed.
    Rerun to get wide page widths right}}

\newcommand{\beginwidepage}[2][1in]{%
  \stepcounter{wpblock}%
  \setlength{\wpmargin}{#1}%
  \ifdim\wpmargin<\wpminmargin
    \setlength{\wpmargin}{\wpminmargin}%   never let text reach the watermark
  \fi
  \settowidth{\wpcharwd}{{\wptabfont 0}}%
  \setlength{\wptblwd}{#2\wpcharwd}%
  %% A width measured on a previous run beats the caller's estimate, but only
  %% upwards: a block may also hold verbatim text, which no table accounts for.
  \@ifundefined{wp@wd@\thewpblock}{}{%
    \setlength{\wpscratch}{\@nameuse{wp@wd@\thewpblock}}%
    \ifdim\wpscratch>\wptblwd
      \setlength{\wptblwd}{\wpscratch}%
    \fi}%
  \global\setlength{\wpcontentwd}{\z@}%
  \global\wpinblocktrue
  \ifdim\wptblwd>\textwidth
    \global\wpwidenedtrue
  \else
    \global\wpwidenedfalse
  \fi
  \ifwpwidened
    \clearpage
    \global\setlength{\wpsaved}{\wppagewidth}%
    \global\setlength{\wpsavedtw}{\textwidth}%
    \global\setlength{\wpsavedosm}{\oddsidemargin}%
    \global\setlength{\wpsavedesm}{\evensidemargin}%
    \global\setlength{\wppagewidth}{\dimexpr\wptblwd+2\wpmargin\relax}%
    \global\setlength{\paperwidth}{\wppagewidth}%
    \global\setlength{\textwidth}{\wptblwd}%
    \global\setlength{\hsize}{\wptblwd}%
    \global\setlength{\linewidth}{\wptblwd}%
    \global\setlength{\columnwidth}{\wptblwd}%
    \global\setlength{\oddsidemargin}{\dimexpr(\paperwidth-\wptblwd)/2-1in\relax}%
    \global\setlength{\evensidemargin}{\oddsidemargin}%
    \typeout{widepage: paper \the\paperwidth\space text \the\textwidth}%
  \else
    \typeout{widepage: table \the\wptblwd\space fits \the\textwidth, no change}%
  \fi
}

\newcommand{\stopwidepage}{%
  \ifdim\wpcontentwd>\z@
    \if@filesw
      \immediate\write\@auxout{%
        \string\wpsetwidth{\thewpblock}{\the\wpcontentwd}}%
    \fi
    \@ifundefined{wp@wd@\thewpblock}{\wp@rerun}{%
      \setlength{\wpscratch}{\@nameuse{wp@wd@\thewpblock}}%
      \ifdim\wpscratch=\wpcontentwd\else\wp@rerun\fi}%
  \fi
  \global\wpinblockfalse
  \ifwpwidened
    \clearpage
    \global\setlength{\wppagewidth}{\wpsaved}%
    \global\setlength{\paperwidth}{\wpsaved}%
    \global\setlength{\textwidth}{\wpsavedtw}%
    \global\setlength{\hsize}{\wpsavedtw}%
    \global\setlength{\linewidth}{\wpsavedtw}%
    \global\setlength{\columnwidth}{\wpsavedtw}%
    \global\setlength{\oddsidemargin}{\wpsavedosm}%
    \global\setlength{\evensidemargin}{\wpsavedesm}%
    \global\wpwidenedfalse
  \fi
}

\makeatother

%% ---- watermark ------------------------------------------------------------
\usepackage{atbegshi}
\usepackage{graphicx}
\usepackage{xcolor}

\newlength{\wpstampinset}   % paper edge -> watermark
\newlength{\wpstampgap}     % watermark -> text block
\setlength{\wpstampinset}{25pt}
\setlength{\wpstampgap}{12pt}

\newcommand{\wpstamp}{%
  \rotatebox{90}{\normalfont\color{red}\fontsize{20pt}{24pt}\selectfont UNVALIDATED CONTENT}%
}

\AtBeginShipout{\AtBeginShipoutUpperLeft{%
  \put(0,-420){%
    \hbox to \paperwidth{\hspace*{\wpstampinset}\wpstamp\hfill\wpstamp\hspace*{\wpstampinset}}%
  }%
}}

\newcommand{\alertwarning}[1]{\par\vspace{4pt}\noindent\fcolorbox{yellow!70!black}{yellow!20}{\parbox{\dimexpr\linewidth-2\fboxsep-2\fboxrule}{\color{yellow!60!black}#1}}\par\vspace{4pt}}

%% ---- long cells -----------------------------------------------------------
%% A listing cell over the length limit is cut and ends in a mark that links to a note, below the table, with the
%% whole text. The note links back. The mark is a single PDF push button: /A follows the link and /TU, the whole
%% text, is its tooltip. A plain link under a tooltip button would not work, the button on top takes the click.
%% Written for xdvipdfmx (xelatex).
\newsavebox{\longcellbox}

%% {characters per line}{text}: a cell that wraps at that many characters of the current, monospace, font
\newcommand{\longcell}[2]{\parbox[t]{#1\fontcharwd\font`0}{\raggedright #2}}

%% Link target. Viewers put it at or near the top of the window, so it is raised three lines to show the line it is
%% on, with some of the one above. \normalbaselineskip, as \baselineskip is 0pt in tables.
\newcommand{\longcellanchor}[1]{\raisebox{3\normalbaselineskip}[0pt][0pt]{\hypertarget{#1}{}}}

%% This is the reason we cannot simply use hyperref and we have to use Forms.
%% We can move to hyperref if we drop the tooltip and the notes
%% {note anchor}{mark}{whole text as a PDF string}
\newcommand{\longcellmark}[3]{%
  \sbox{\longcellbox}{\textcolor{blue}{#2}}%
  \special{pdf:ann width \the\wd\longcellbox\space height \the\ht\longcellbox\space depth \the\dp\longcellbox\space
    << /Subtype /Widget /FT /Btn /Ff 65536 /T (#1) /TU #3 /A << /S /GoTo /D (#1) >> /F 4 /H /I /BS << /W 0 >> >>}%
  \usebox{\longcellbox}}

\newenvironment{longcellnotes}{\par\raggedright}{\par}

%% {note anchor}{cell anchor}{mark}{whole text, escaped}
\newcommand{\longcellnote}[4]{%
  \par\noindent
  \sbox{\longcellbox}{#3\ }%
  \hangindent\wd\longcellbox
  \longcellanchor{#1}%
  \hyperlink{#2}{\textcolor{blue}{#3}}\ #4\par}

%% ---- table of contents rerun ------------------------------------------------
%% Page breaks can still move on the second LaTeX pass (wide pages, a table of contents that outgrows the title
%% page) and LaTeX does not warn when the .toc changes, so tinytex stops one pass early and prints stale page
%% numbers. Ask for a rerun in words tinytex recognizes. The .toc must be closed before it is checksummed.
\usepackage{rerunfilecheck}
\makeatletter
\RerunFileCheck{\jobname.toc}{%
  \@ifundefined{tf@toc}{}{\immediate\closeout\tf@toc}%
}{Rerun to get the table of contents right}
\makeatother
